\documentclass[a4paper]{article}     
\usepackage{amsfonts}
\usepackage{amsmath}
\usepackage{amssymb}
\usepackage{graphicx}
\usepackage{color}

\newcommand{\Q}{\mathbb{Q}}
\newcommand{\R}{\mathbb{R}}
\newcommand{\llim}{\lim\limits}
\newcommand{\dint}{\displaystyle\int}

\begin{document}

\noindent \large Auteur: Hina Rana \\
\noindent \large Studierichting: Informatica\\
\noindent \large Oefeningenreeks 2E nr. 12.7\\

\medskip

\normalsize

$\lim\limits_{x \to \infty} \sqrt{25x^2-6}-5x$\\

$\lim\limits_{x \to -\infty} \sqrt{25x^2-6}-5x = + \infty + \infty = + \infty$\\

$\lim\limits_{x \to +\infty} \sqrt{25x^2-6}-5x = + \infty - \infty$\\

$=\lim\limits_{x \to +\infty} \dfrac{(\sqrt{25x^2-6}-5x)(\sqrt{25x^2-6}+5x)}{\sqrt{25x^2-6}+5x}$\\

$=\lim\limits_{x \to +\infty} \dfrac{25x^2-6-(-5x)^2}{\sqrt{25x^2-6}+5x} = \dfrac{25x^2-6-25x^2}{\sqrt{25x^2-6}+5x}$\\

$=\lim\limits_{x \to +\infty} \dfrac{-6}{\sqrt{25x^2-6}+5x} = 0$\\







\begin{thebibliography}{999}
\end{thebibliography}
\end{document} 
